\documentclass{otbook}

\usepackage{otsymbols}

\title{Nomenclature and Symbol Management}
\author{O. T. Author}
\date{\today}

% --- Symbol Definitions ---
% Notice how both use \rho but have different keys
\otdefinesymbol{mass-density}{\rho}{Density}{\unit{kg.m^{-3}}}{The mass per unit volume of a substance.}
\otdefinesymbol{resistivity}{\rho}{Resistivity}{\unit{\ohm.m}}{The inherent property of a material to resist electrical current.}

\otdefinesymbol{pressure}{P}{Pressure}{\unit{Pa}}{The force applied perpendicular to the surface of an object per unit area.}
\otdefinesymbol{power}{P}{Power}{\unit{W}}{The rate at which work is done or energy is transferred.}

% --- Acronym Definitions ---
\otdefineacronym{ota}{OTA}{Original Tech Author}{The primary maintainer of the LaTeX Lab suite.}

\begin{document}

\makeotbooktitle

\chapter{The Collision Problem}

In scientific writing, the same Greek symbol is often used for different physical quantities. For example, \sym{mass-density} represents mass density, while \sym{resistivity} represents electrical resistivity. Both appear as $\rho$, but they are distinct concepts.

Similarly, we might use $P$ for \sym{pressure} and \sym{power}. By using unique keys, we ensure that the nomenclature table tracks them separately.

We can also use acronyms like \acr{ota}. On first use, it expands. On subsequent use, like \acr{ota}, it shows only the short form.

\chapter{Nomenclature Table}

\printotnomenclature[Table of Symbols and Abbreviations]

\end{document}
